\documentclass[twocolumn,10pt]{article}
\usepackage{iftex}
\ifPDFTeX
    \usepackage[utf8]{inputenc}
    \usepackage[T1]{fontenc}
\else
    \usepackage{fontspec}
    \setmainfont{TeX Gyre Termes}
    \setsansfont{TeX Gyre Heros}
    \setmonofont{TeX Gyre Cursor}
\fi
\usepackage{kotex}
\ifPDFTeX\else
    \setmainhangulfont[AutoFakeSlant=0.15]{Noto Serif CJK KR}
    \setsanshangulfont[AutoFakeSlant=0.15]{Noto Sans CJK KR}
\fi
\usepackage{amsmath, amssymb, amsfonts}
\usepackage{booktabs}
\usepackage{array}
\usepackage{tabularx}
\usepackage{graphicx}
\usepackage[table,dvipsnames]{xcolor}
\usepackage{tikz}
\usetikzlibrary{arrows.meta,positioning,calc,fit,shapes.geometric,decorations.pathreplacing,backgrounds}
\usepackage{pgfplots}
\pgfplotsset{compat=1.18}
\usepackage{caption}
\usepackage{float}
\usepackage{cite}
\usepackage{geometry}
\usepackage{microtype}
\usepackage{ragged2e}
\usepackage{hyperref}
\geometry{a4paper, margin=1.8cm}
\setlength{\columnsep}{0.7cm}
\setlength{\parskip}{0.18em}
\emergencystretch=2em

\definecolor{AIInk}{HTML}{1B2633}
\definecolor{CloudBlue}{HTML}{2868C7}
\definecolor{EdgeGreen}{HTML}{1E8A5A}
\definecolor{WarmGold}{HTML}{C88719}
\definecolor{AlertRed}{HTML}{C84A4A}
\definecolor{PanelGray}{HTML}{F2F5F7}
\definecolor{LineGray}{HTML}{6D7B88}

\hypersetup{
    colorlinks=true,
    linkcolor=CloudBlue,
    citecolor=EdgeGreen,
    urlcolor=CloudBlue
}
\captionsetup{font=small,labelfont=bf}
\tikzset{
    stage/.style={draw=AIInk!55, fill=#1, rounded corners=2pt, very thick, align=center, inner sep=5pt, font=\sffamily\small},
    chip/.style={draw=AIInk!55, fill=#1, rounded corners=1pt, line width=0.7pt, align=center, inner sep=4pt, font=\sffamily\scriptsize},
    flow/.style={-{Stealth[length=2.2mm]}, very thick, draw=AIInk!78},
    softflow/.style={-{Stealth[length=2mm]}, thick, draw=LineGray!85},
    metric/.style={draw=AIInk!45, fill=white, rounded corners=1.5pt, align=center, inner sep=3pt, font=\sffamily\scriptsize}
}
\renewcommand{\figurename}{Fig.}
\renewcommand{\tablename}{Table}

\begin{document}

\twocolumn[
\begin{@twocolumnfalse}
\begin{center}
{\LARGE\bfseries 온디바이스 소형 언어모델의 이해: 클라우드 AI에서 개인 기기 AI로\par}
\vspace{5pt}
{\large Understanding On-device Small Language Models: From Cloud AI to Personal Device AI\par}
\vspace{8pt}
{\normalsize 김현우\\Hyun Woo Kim\par}
\vspace{6pt}
{\small 2026년 4월 22일\par}
\end{center}

\vspace{6pt}
\begin{center}
\begin{minipage}{0.86\textwidth}
\small
\noindent\textbf{Abstract}\quad
This paper explains the recent move from cloud-based large language models to on-device small language models in an easy-to-read way. Cloud AI is powerful, but it can be expensive, slow when the network is poor, and sensitive for private data. On-device AI runs a smaller model directly on a phone, laptop, car, or wearable device. To make this practical, the model must use less memory, the chip must move data efficiently, and the system must choose the right place to run each task. The paper also introduces Google DeepMind's AlphaEvolve as an example of AI that searches for better algorithms by writing, testing, and improving code. This paper focuses on simple ideas such as quantization, cache compression, mobile NPU scheduling, privacy protection, and personal-device collaboration.
\vspace{4pt}

\noindent\textbf{Key Words}: On-device AI, Small Language Model, Quantization, Mobile NPU, Privacy, Edge AI
\end{minipage}
\end{center}
\vspace{14pt}
\end{@twocolumnfalse}
]

\RaggedRight

\section{서론}
최근 생성형 인공지능은 대부분 클라우드 서버에서 동작했다. 사용자가 질문을 보내면, 서버의 큰 GPU가 답을 만들고 다시 사용자에게 전송하는 방식이다. 이 구조는 성능이 좋지만 몇 가지 단점도 있다. 네트워크가 느리면 답이 늦게 오고, 서버 비용이 많이 들며, 개인적인 내용이 외부 서버로 전송될 수 있다.

이 문제를 줄이기 위해 온디바이스 AI가 주목받고 있다. 온디바이스 AI는 스마트폰, 노트북, 자동차, 웨어러블 기기 안에서 모델을 직접 실행하는 방식이다. 최근 Llama 3.2, Phi 계열, Gemma 계열처럼 비교적 작은 언어모델이 공개되면서, 작은 기기에서도 쓸 만한 AI 기능을 실행할 가능성이 커졌다 \cite{llama32_2024,phi4_2025,gemma3_2025}.

하지만 모델을 작게 만드는 것만으로는 충분하지 않다. 실제 기기에는 배터리, 발열, 메모리, 보안, 사용자의 체감 속도 같은 제한이 있기 때문이다. 따라서 온디바이스 AI를 이해하려면 ``모델 압축''뿐 아니라 ``기기에서 어떻게 실행할지''까지 함께 봐야 한다. Fig. \ref{fig:ko_bigpicture}는 이 흐름을 한눈에 정리한 그림이다.

\begin{figure}[H]
\centering
\resizebox{\columnwidth}{!}{%
\begin{tikzpicture}[node distance=0.52cm]
    \node[stage=CloudBlue!12, text width=2.1cm, minimum height=0.85cm] (cloud)
        {\textbf{클라우드 AI}\\\footnotesize 강하지만\\멀리 있음};
    \node[stage=PanelGray, text width=2.05cm, minimum height=0.85cm, right=of cloud] (need)
        {\textbf{필요}\\\footnotesize 빠르고\\개인적인 AI};
    \node[stage=EdgeGreen!12, text width=2.1cm, minimum height=0.85cm, right=of need] (device)
        {\textbf{온디바이스 AI}\\\footnotesize 작지만\\가까이 있음};

    \draw[flow] (cloud) -- (need);
    \draw[flow] (need) -- (device);

    \node[chip=WarmGold!16, text width=1.75cm, below=0.42cm of cloud] (c1) {모델 줄이기};
    \node[chip=WarmGold!16, text width=1.75cm, below=0.42cm of need] (c2) {데이터 줄이기};
    \node[chip=WarmGold!16, text width=1.75cm, below=0.42cm of device] (c3) {기기 상태 보기};
    \draw[softflow] (c1) -- (c2);
    \draw[softflow] (c2) -- (c3);
\end{tikzpicture}%
}
\caption{클라우드 AI에서 온디바이스 AI로 이동하는 이유. 핵심은 작은 기기에서 빠르고 오래 쓸 수 있게 만드는 것이다.}
\label{fig:ko_bigpicture}
\end{figure}

\section{새로운 흐름: 알고리즘을 개선하는 AI}
지금까지의 생성형 AI는 주로 글을 쓰거나 질문에 답하는 도구로 이해되었다. 그런데 Google DeepMind의 AlphaEvolve는 조금 다른 방향을 보여준다. AlphaEvolve는 답변 문장을 만드는 데서 끝나지 않고, 코드로 표현된 알고리즘을 직접 바꾸고, 실행해 보고, 점수를 비교하면서 더 나은 알고리즘을 찾는다 \cite{deepmind_alphaevolve_2025}.

직관적으로 말하면 다음과 같다.
\begin{center}
\small
\textbf{아이디어 만들기} $\rightarrow$ \textbf{코드 실행} $\rightarrow$ \textbf{점수 매기기} $\rightarrow$ \textbf{좋은 아이디어 다시 사용}
\end{center}
이 과정이 반복되면, 사람이 처음에 넣은 단순한 코드보다 더 좋은 코드가 나올 수 있다. 여기서 중요한 점은 ``AI가 말했다''가 아니라 ``테스트를 통과했고 점수가 좋아졌다''이다.

\begin{figure}[H]
\centering
\resizebox{\columnwidth}{!}{%
\begin{tikzpicture}[node distance=0.5cm and 0.55cm]
    \node[stage=PanelGray, text width=1.65cm] (human) {사람\\\footnotesize 문제와 점수 정의};
    \node[stage=CloudBlue!10, text width=1.65cm, right=of human] (llm) {LLM\\\footnotesize 코드 수정 제안};
    \node[stage=WarmGold!16, text width=1.65cm, right=of llm] (test) {평가기\\\footnotesize 실행/검증};
    \node[stage=EdgeGreen!10, text width=1.65cm, right=of test] (db) {저장소\\\footnotesize 좋은 코드 보관};

    \draw[flow] (human) -- (llm);
    \draw[flow] (llm) -- (test);
    \draw[flow] (test) -- (db);
    \draw[softflow] (db) to[bend left=22] node[below,font=\sffamily\scriptsize] {다음 세대에 재사용} (llm);

    \node[chip=CloudBlue!10, text width=1.55cm, below=0.45cm of llm] (a) {스케줄링};
    \node[chip=EdgeGreen!10, text width=1.55cm, right=0.28cm of a] (b) {커널 최적화};
    \node[chip=AlertRed!8, text width=1.55cm, right=0.28cm of b] (c) {칩 회로};
\end{tikzpicture}%
}
\caption{AlphaEvolve의 핵심 흐름. 사람이 평가 기준을 만들고, AI는 코드를 바꿔 가며 더 높은 점수를 찾는다.}
\label{fig:ko_alphaevolve}
\end{figure}

Google DeepMind는 AlphaEvolve가 여러 실제 문제에서 성과를 냈다고 보고했다. 예를 들어 데이터센터 스케줄링을 개선해 전 세계 컴퓨팅 자원의 약 0.7\%를 회수했고, TPU 회로의 불필요한 부분을 줄이는 Verilog 수준의 수정을 제안했으며, Gemini 학습에 쓰이는 행렬 곱셈 커널을 더 빠르게 만들었다 \cite{deepmind_alphaevolve_2025}. 또한 4$\times$4 복소수 행렬 곱셈에서 48번의 스칼라 곱셈만 사용하는 방법을 찾는 등 수학적 알고리즘 발견에도 활용되었다.

이 내용이 온디바이스 AI 논문에 중요한 이유는 간단하다. 작은 모델을 빠르게 돌리려면 양자화만 잘하면 되는 것이 아니다. 어떤 커널을 쓸지, 메모리를 어떻게 배치할지, NPU와 CPU 중 어디에 일을 맡길지 같은 선택도 매우 중요하다. AlphaEvolve식 반복 탐색은 이런 선택을 자동으로 실험하고 개선할 수 있는 가능성을 보여준다. 따라서 앞으로의 하드웨어-소프트웨어 공동 설계는 사람이 규칙을 모두 손으로 짜는 방식에서, 사람이 목표를 정하고 AI가 후보 설계를 많이 시험하는 방식으로 확장될 수 있다.

\section{핵심 개념: 모델을 가볍게 만들기}
\subsection{양자화}
언어모델은 매우 많은 숫자, 즉 가중치로 이루어져 있다. 보통 정확도를 높이기 위해 이 숫자를 자세하게 저장하지만, 모바일 기기에서는 그만큼 메모리를 많이 사용한다. 양자화는 이 숫자를 더 적은 비트로 표현하는 방법이다. 예를 들어 16비트 숫자를 4비트나 2비트로 줄이면, 저장 공간과 메모리 이동량을 크게 줄일 수 있다 \cite{frantar2022gptq,lin2023awq}.

가장 직관적인 예는 가중치를 세 값으로만 표현하는 방식이다. 복잡한 소수 대신 $-1$, $0$, $+1$ 중 하나로 바꾸면, 곱셈을 덧셈이나 생략으로 바꿀 수 있다.
\begin{equation}
w_q \in \{-1,\ 0,\ +1\}.
\label{eq:ko_simple_quant}
\end{equation}
식 \ref{eq:ko_simple_quant}은 매우 단순하지만 의미가 크다. $0$이 된 값은 계산하지 않아도 되고, $+1$과 $-1$은 곱셈기 없이 더하거나 빼는 방식으로 처리할 수 있다. Fig. \ref{fig:ko_ternary}는 이 과정을 숫자를 세 바구니로 나누는 모습으로 표현한 것이다.

\begin{figure}[t]
\centering
\begin{tikzpicture}[x=0.86cm,y=0.8cm,font=\sffamily\scriptsize]
    \draw[->, thick, draw=AIInk] (-3.2,0) -- (3.25,0) node[right] {원래 가중치};
    \draw[densely dashed, draw=LineGray] (-1.25,-0.15) -- (-1.25,1.55);
    \draw[densely dashed, draw=LineGray] (1.25,-0.15) -- (1.25,1.55);

    \node[stage=AlertRed!10, text width=1.55cm, minimum height=0.85cm] at (-2.25,0.9) {$-1$\\\footnotesize 빼기};
    \node[stage=WarmGold!16, text width=1.55cm, minimum height=0.85cm] at (0,0.9) {$0$\\\footnotesize 생략};
    \node[stage=EdgeGreen!12, text width=1.55cm, minimum height=0.85cm] at (2.25,0.9) {$+1$\\\footnotesize 더하기};

    \node[below] at (-1.25,-0.15) {작음};
    \node[below] at (1.25,-0.15) {큼};
    \draw[decorate,decoration={brace,mirror,amplitude=4pt},draw=LineGray]
        (-1.2,-0.6) -- (1.2,-0.6) node[midway,below=5pt] {작은 값은 0으로 정리};
\end{tikzpicture}
\caption{양자화의 직관적 의미. 복잡한 숫자를 몇 개의 대표값으로 바꾸면 메모리와 계산량이 줄어든다.}
\label{fig:ko_ternary}
\end{figure}

\subsection{KV-cache 압축}
언어모델이 긴 글을 읽거나 긴 대화를 이어갈 때는 이전 문맥을 기억해야 한다. 이때 사용하는 임시 기억 공간을 KV-cache라고 부른다. 문맥이 길어질수록 이 공간도 함께 커지기 때문에, 작은 기기에서는 큰 부담이 된다 \cite{kvquant2024}.

이 관계는 다음처럼 간단히 볼 수 있다.
\begin{equation}
\text{KV-cache 크기} \propto \text{문맥 길이} \times \text{저장 비트 수}.
\label{eq:ko_kv_simple}
\end{equation}
즉, 더 긴 대화를 처리하려면 더 많은 메모리가 필요하다. 그래서 온디바이스 AI에서는 모델 가중치뿐 아니라 KV-cache도 함께 압축해야 한다.

\section{간단한 비교 실험}
Table \ref{tab:ko_results}는 같은 소형 언어모델을 서로 다른 정밀도로 저장했을 때의 예시 비교이다. 숫자는 실제 기기와 모델에 따라 달라질 수 있지만, 전체 흐름을 이해하는 데 도움이 된다. 비트 수가 줄어들수록 메모리 사용량이 줄고, 같은 전력으로 더 많은 토큰을 처리할 수 있다.

\begin{table*}[t]
\centering
\caption{정밀도에 따른 온디바이스 추론 예시 비교}
\label{tab:ko_results}
\begin{tabular}{lcccc}
\toprule
\textbf{정밀도} & \textbf{메모리 사용량} & \textbf{속도} & \textbf{토큰당 에너지} & \textbf{해석} \\
\midrule
FP16 & 7.6 GB & 12.5 tok/s & 120.4 mJ & 정확하지만 무겁다 \\
INT8 & 3.9 GB & 38.2 tok/s & 42.1 mJ & 기본적인 경량화 \\
INT4 & 2.1 GB & 85.5 tok/s & 18.2 mJ & 모바일에 더 적합 \\
\rowcolor{EdgeGreen!10}
1.58-bit & 1.2 GB & 142.0 tok/s & 6.4 mJ & 매우 작고 빠르다 \\
\bottomrule
\end{tabular}
\end{table*}

\begin{figure*}[t]
\centering
\begin{tikzpicture}
\begin{axis}[
    width=0.47\textwidth,
    height=4.8cm,
    ybar,
    bar width=12pt,
    ymin=0,
    ylabel={속도 (tokens/s)},
    symbolic x coords={FP16,INT8,INT4,Ternary},
    xtick=data,
    xticklabels={FP16,INT8,INT4,1.58-bit},
    nodes near coords,
    nodes near coords style={font=\sffamily\scriptsize},
    tick label style={font=\sffamily\scriptsize},
    label style={font=\sffamily\small},
    ymajorgrids,
    grid style={LineGray!18},
    axis line style={draw=AIInk!70},
    enlarge x limits=0.16
]
\addplot[fill=CloudBlue!65,draw=CloudBlue!90] coordinates {(FP16,12.5) (INT8,38.2) (INT4,85.5) (Ternary,142.0)};
\end{axis}
\end{tikzpicture}
\hfill
\begin{tikzpicture}
\begin{axis}[
    width=0.47\textwidth,
    height=4.8cm,
    ybar,
    bar width=12pt,
    ymin=0,
    ylabel={에너지 (mJ/token)},
    symbolic x coords={FP16,INT8,INT4,Ternary},
    xtick=data,
    xticklabels={FP16,INT8,INT4,1.58-bit},
    nodes near coords,
    nodes near coords style={font=\sffamily\scriptsize},
    tick label style={font=\sffamily\scriptsize},
    label style={font=\sffamily\small},
    ymajorgrids,
    grid style={LineGray!18},
    axis line style={draw=AIInk!70},
    enlarge x limits=0.16
]
\addplot[fill=AlertRed!62,draw=AlertRed!90] coordinates {(FP16,120.4) (INT8,42.1) (INT4,18.2) (Ternary,6.4)};
\end{axis}
\end{tikzpicture}
\caption{정밀도별 속도와 에너지 비교. 왼쪽 그래프는 높을수록 좋고, 오른쪽 그래프는 낮을수록 좋다.}
\label{fig:ko_perf}
\end{figure*}

Fig. \ref{fig:ko_perf}에서 중요한 점은 단순히 ``빠르다''가 아니다. 모바일 기기에서는 배터리와 발열이 함께 중요하다. 같은 답을 만들 때 에너지를 적게 쓰면 배터리가 오래가고, 기기가 뜨거워지는 문제도 줄어든다.

\section{기기에서 잘 돌리기 위한 시스템 설계}
\subsection{모델만 줄이면 끝이 아닌 이유}
소형 언어모델을 실제 제품에 넣으려면 모델 파일만 작게 만들면 안 된다. 모델이 실행되는 동안 데이터가 메모리와 NPU 사이를 계속 이동하고, 이 이동이 전력과 시간을 많이 사용한다. 그래서 하드웨어와 소프트웨어를 함께 설계해야 한다 \cite{roofline2009,ghose2019survey}.

예를 들어 1.58-bit 모델을 사용한다면, NPU도 그 형식에 맞게 동작하는 것이 좋다. 일반적인 곱셈기를 그대로 쓰는 대신, $-1$, $0$, $+1$에 맞춘 간단한 덧셈 구조를 쓰면 더 효율적이다. Fig. \ref{fig:ko_codesign}은 이 개념을 단순화한 그림이다.

\begin{figure*}[t]
\centering
\begin{tikzpicture}[node distance=0.65cm and 0.9cm]
    \node[stage=PanelGray, text width=2.55cm, minimum height=1.0cm] (app) {앱 요청\\\footnotesize 질문, 음성, 사진};
    \node[stage=WarmGold!16, text width=2.75cm, minimum height=1.0cm, right=of app] (scheduler) {실행 관리자\\\footnotesize 빠른지, 안전한지\\배터리는 충분한지 확인};
    \node[stage=CloudBlue!10, text width=2.75cm, minimum height=1.0cm, right=of scheduler] (npu) {모바일 NPU\\\footnotesize 작은 모델 계산};
    \node[stage=EdgeGreen!10, text width=2.75cm, minimum height=1.0cm, right=of npu] (memory) {메모리 관리\\\footnotesize 필요한 데이터만 이동};

    \draw[flow] (app) -- (scheduler);
    \draw[flow] (scheduler) -- (npu);
    \draw[flow] (npu) -- (memory);
    \draw[softflow] (memory) to[bend left=12] node[below,font=\sffamily\scriptsize] {결과와 상태 정보} (scheduler);

    \node[metric, below=0.45cm of scheduler] (m1) {발열};
    \node[metric, right=0.35cm of m1] (m2) {배터리};
    \node[metric, right=0.35cm of m2] (m3) {개인정보};
    \node[metric, right=0.35cm of m3] (m4) {응답 시간};
\end{tikzpicture}
\caption{온디바이스 AI 실행 구조. 실행 관리자는 단순히 가장 빠른 방법이 아니라, 현재 기기 상태에 맞는 방법을 선택한다.}
\label{fig:ko_codesign}
\end{figure*}

\subsection{실행 위치 선택}
항상 기기 안에서만 실행하는 것이 정답은 아니다. 매우 어려운 질문이나 큰 파일 분석은 클라우드가 더 적합할 수 있다. 반대로 개인정보가 많은 내용이나 네트워크가 불안정한 상황은 기기 안에서 처리하는 것이 좋다. 이를 간단히 점수로 표현하면 다음과 같다.
\begin{equation}
\text{선택 점수} = \text{지연 시간} + \text{에너지} + \text{개인정보 위험}.
\label{eq:ko_simple_score}
\end{equation}
점수가 낮은 쪽을 선택하면 된다. 실제 시스템에서는 여기에 모델 크기, 배터리 상태, 사용자의 설정 등이 추가된다.

\begin{figure}[t]
\centering
\resizebox{\columnwidth}{!}{%
\begin{tikzpicture}[node distance=0.55cm and 0.65cm]
    \node[stage=PanelGray, text width=1.75cm] (query) {사용자 요청};
    \node[draw=WarmGold!80!black, fill=WarmGold!14, diamond, aspect=1.75, align=center, inner sep=1.5pt, font=\sffamily\scriptsize, right=of query] (privacy) {민감한\\정보?};
    \node[chip=EdgeGreen!12, text width=1.6cm, above right=0.05cm and 0.65cm of privacy] (local) {기기에서\\처리};
    \node[draw=WarmGold!80!black, fill=WarmGold!14, diamond, aspect=1.75, align=center, inner sep=1.5pt, font=\sffamily\scriptsize, below right=0.05cm and 0.65cm of privacy] (hard) {너무\\어려움?};
    \node[chip=CloudBlue!12, text width=1.6cm, right=0.65cm of hard] (cloud) {클라우드\\도움};
    \node[chip=EdgeGreen!12, text width=1.6cm, above=0.75cm of cloud] (fast) {작은 모델로\\빠르게 처리};

    \draw[flow] (query) -- (privacy);
    \draw[flow] (privacy) -- node[above,font=\sffamily\scriptsize] {예} (local);
    \draw[softflow] (privacy) -- node[left,font=\sffamily\scriptsize] {아니오} (hard);
    \draw[flow] (hard) -- node[below,font=\sffamily\scriptsize] {예} (cloud);
    \draw[softflow] (hard) -- node[right,font=\sffamily\scriptsize] {아니오} (fast);
\end{tikzpicture}%
}
\caption{온디바이스와 클라우드 선택의 직관적 기준. 개인정보와 난이도를 먼저 확인하면 이해하기 쉽다.}
\label{fig:ko_decision}
\end{figure}

\section{개인정보와 보안}
온디바이스 AI의 큰 장점은 사용자의 데이터가 기기 밖으로 나가지 않아도 된다는 점이다. 예를 들어 개인 일정, 문자, 사진 설명, 음성 메모처럼 민감한 정보는 기기 안에서 처리하는 편이 안전하다. 그러나 이것만으로 모든 보안 문제가 해결되는 것은 아니다. 기기 내부에서도 악성 앱이나 운영체제 취약점이 있을 수 있기 때문이다.

그래서 최근에는 보안 실행 환경(TEE)이나 private cloud 방식이 함께 사용된다 \cite{arm_cca_2024,apple_pcc_2024}. TEE는 중요한 계산을 별도의 보호된 공간에서 실행하는 방식이다. Private cloud는 꼭 서버 도움이 필요할 때도 사용자의 데이터를 강하게 보호하도록 설계된 서버 구조이다.

\begin{figure}[t]
\centering
\resizebox{\columnwidth}{!}{%
\begin{tikzpicture}[node distance=0.45cm and 0.55cm]
    \node[stage=PanelGray, text width=1.45cm] (data) {개인 데이터};
    \node[stage=EdgeGreen!10, text width=1.45cm, right=of data] (local) {기기 안\\처리};
    \node[stage=WarmGold!16, text width=1.45cm, right=of local] (safe) {보호 영역\\TEE};
    \node[stage=CloudBlue!10, text width=1.45cm, right=of safe] (pcc) {필요시\\사설 클라우드};
    \node[stage=PanelGray, text width=1.45cm, right=of pcc] (answer) {답변};

    \draw[flow] (data) -- (local);
    \draw[flow] (local) -- (safe);
    \draw[softflow] (safe) -- node[above,font=\sffamily\scriptsize] {선택적} (pcc);
    \draw[flow] (pcc) -- (answer);
\end{tikzpicture}%
}
\caption{개인정보 보호 관점의 실행 흐름. 민감한 처리는 가능한 한 기기 안에서 수행하고, 필요한 경우에만 보호된 서버를 사용한다.}
\label{fig:ko_privacy}
\end{figure}

\section{멀티모달 AI와 개인 기기 협업}
앞으로의 온디바이스 AI는 텍스트만 다루지 않는다. 음성, 사진, 카메라 영상, 센서 데이터까지 함께 처리해야 한다. 이때 모든 기능을 하나의 큰 모델이 담당하면 너무 무겁다. 대신 필요한 부분만 켜는 방식이 유리하다. 예를 들어 사진 질문에는 vision 모듈을, 음성 대화에는 audio 모듈을 더 많이 사용하는 식이다 \cite{minicpm_v_2025,moshi2024}.

또 다른 방향은 여러 개인 기기가 서로 도와주는 것이다. 스마트폰은 항상 가까이에 있지만 작고, 노트북은 더 강하지만 항상 들고 다니지는 않는다. 태블릿이나 PC까지 함께 사용하면, 한 기기만 사용할 때보다 더 큰 작업을 처리할 수 있다. Fig. \ref{fig:ko_personal_cloud}는 이런 개인 기기 협업 구조를 보여준다.

\begin{figure}[t]
\centering
\resizebox{\columnwidth}{!}{%
\begin{tikzpicture}[node distance=0.58cm and 0.6cm]
    \node[stage=CloudBlue!10, text width=1.55cm] (phone) {스마트폰\\\footnotesize 빠른 입력};
    \node[stage=EdgeGreen!10, text width=1.55cm, right=of phone] (tablet) {태블릿\\\footnotesize 화면/문서};
    \node[stage=WarmGold!16, text width=1.55cm, right=of tablet] (pc) {노트북\\\footnotesize 무거운 계산};
    \node[metric, below=0.45cm of tablet] (net) {개인 네트워크로 연결};

    \draw[flow] (phone) -- (tablet);
    \draw[flow] (tablet) -- (pc);
    \draw[softflow] (net) -- (phone);
    \draw[softflow] (net) -- (tablet);
    \draw[softflow] (net) -- (pc);
\end{tikzpicture}%
}
\caption{개인 기기 협업 구조. 여러 기기가 역할을 나누면 작은 기기의 한계를 줄일 수 있다.}
\label{fig:ko_personal_cloud}
\end{figure}

개인화도 중요한 주제이다. 사용자의 말투, 자주 쓰는 앱, 일정 패턴을 반영하면 AI가 더 편리해진다. 이때 전체 모델을 다시 학습하는 것은 너무 비싸다. 대신 작은 추가 부품만 학습하는 PEFT 방식이 많이 사용된다 \cite{dora2024,minions2025}. 이를 간단히 쓰면 다음과 같다.
\begin{equation}
\text{개인화 모델} = \text{기본 모델} + \text{작은 개인화 부품}.
\label{eq:ko_peft_simple}
\end{equation}
이 방식은 기본 모델은 그대로 두고, 사용자에게 필요한 작은 부분만 업데이트한다는 점에서 온디바이스 환경과 잘 맞는다.

\section{결론}
본 논문은 클라우드 AI에서 온디바이스 AI로 이동하는 흐름을 학부생 수준에서 이해할 수 있도록 정리하였다. 클라우드 AI는 강력하지만 비용, 지연 시간, 개인정보 문제가 있다. 온디바이스 AI는 이러한 문제를 줄일 수 있지만, 작은 기기 안에서 실행해야 하므로 메모리, 배터리, 발열, 보안을 함께 고려해야 한다.

핵심 기술은 크게 세 가지로 볼 수 있다. 첫째, 양자화와 KV-cache 압축으로 모델과 중간 데이터를 가볍게 만든다. 둘째, 모바일 NPU와 메모리 구조를 모델 형식에 맞게 설계한다. 셋째, 개인정보와 기기 상태를 고려해 기기 실행과 클라우드 실행을 적절히 선택한다. 결국 온디바이스 AI의 목표는 단순히 가장 큰 모델을 쓰는 것이 아니라, 사용자의 기기에서 빠르고 안전하며 오래 사용할 수 있는 AI를 만드는 것이다.

\bibliographystyle{plain}
\bibliography{references}

\end{document}
